%==============================================================================
%  CAPÍTULO XVI
%==============================================================================
\chapter{Manual de tramitación: reorganización simplificada de MIPEs}
\label{cap:tram-reorg-simplificada}

\fichaProcedimiento
  {Micro y pequeña empresa (hasta \uf{25.000} de ventas anuales).}
  {Juzgado civil del domicilio del deudor, con veedor de Categoría B.}
  {De tres a seis meses hasta el acuerdo firme.}
  {Reducido: sin auditoría externa obligatoria y con aranceles del veedor acotados.}
  {Acuerdo de reorganización simplificado, con posibilidad de reintento si se rechaza.}

El desarrollo dogmático se encuentra en el capítulo
\ref{cap:reorganizacion-simplificada}.

%==============================================================================
\bloqueEncargo
%==============================================================================

Este procedimiento fue diseñado para un deudor que carece de estructura administrativa.
Esa misma característica define el encargo: buena parte del trabajo que en la
reorganización ordinaria realiza el área de administración del cliente, aquí lo realiza
el abogado.

\begin{foco}[La declaración jurada reemplaza a la auditoría, no a la verdad]
La sustitución de la auditoría externa por declaraciones juradas suscritas por el propio
deudor abarata el procedimiento y traslada el riesgo. Lo que un auditor habría detectado
antes de presentar, ahora lo detectará un acreedor durante la tramitación. La
verificación documental previa del capítulo \ref{cap:entrevista} deja de ser una buena
práctica y pasa a ser la única salvaguarda.
\end{foco}

%==============================================================================
\bloqueEntrevista
%==============================================================================

\begin{cuestionario}[Elegibilidad y estructura]
  \pregunta{¿Cuáles fueron las ventas de los últimos ejercicios, según la carpeta
  tributaria?}
  {Define la elegibilidad para la vía simplificada. Es el primer punto que un acreedor
  objetará si el deudor está cerca del umbral.}

  \pregunta{¿Quién lleva la contabilidad y con qué grado de actualización?}
  {Determina cuánto trabajo de reconstrucción documental deberá asumirse antes de
  presentar.}

  \pregunta{¿El negocio depende de una persona determinada?}
  {En la MIPE la continuidad suele depender del dueño. Es un dato que la propuesta debe
  abordar explícitamente para ser creíble.}

  \pregunta{¿Existen deudas personales del socio confundidas con las de la empresa?}
  {Situación habitual en la microempresa. Puede exigir dos procedimientos paralelos o
  redefinir el sujeto del concurso.}
\end{cuestionario}

\begin{cuestionario}[Preparación del reintento]
  \pregunta{Si la primera propuesta fuera rechazada, ¿qué concesión adicional está
  dispuesto a ofrecer el deudor?}
  {El mecanismo de reintento opera en un plazo brevísimo. La segunda propuesta debe estar
  pensada antes de que se vote la primera.}

  \pregunta{¿Qué acreedores sumarían más del 50\,\% del pasivo con derecho a voto?}
  {Es el apoyo que el reintento exige. Identificarlos con antelación es lo que hace
  utilizable el mecanismo.}
\end{cuestionario}

%==============================================================================
\bloqueDocumental
%==============================================================================

\begin{checklist}
  \item Declaración jurada simple de activos y pasivos suscrita por el deudor.
  \item Carpeta tributaria completa que acredite el umbral de ventas.
  \item Nómina de acreedores con montos, garantías y preferencias.
  \item Propuesta de acuerdo de reorganización simplificada (modelo en el capítulo
  \ref{cap:reorganizacion-simplificada}).
  \item Antecedentes de respaldo de cada partida declarada, aunque la ley no los exija:
  constituyen la defensa frente a objeciones.
\end{checklist}

%==============================================================================
\bloqueFocos
%==============================================================================

\subsection{El plazo de protección es más breve}

El período de protección simplificado es sustancialmente más corto que el ordinario. La
consecuencia práctica es que no hay margen para negociar durante la tramitación: la
negociación debe estar cerrada, o muy avanzada, al presentar.

\subsection{El mecanismo de reintento}

El derecho de proponer un nuevo acuerdo tras el rechazo es la principal ventaja del
procedimiento simplificado, y la más desaprovechada. Su plazo es exiguo y exige acreditar
apoyo cualificado. Sólo puede usarse si se preparó de antemano.

\begin{foco}[Preparar la segunda propuesta antes de la primera votación]
Redactar la propuesta de reintento y obtener las adhesiones necesarias antes de la junta
convierte un rechazo en una segunda oportunidad. No hacerlo convierte el rechazo en
liquidación.
\end{foco}

\subsection{El veedor de Categoría B}

La nominación recae en veedores especializados con aranceles acotados. La contrapartida
es que suelen manejar mayor volumen de causas. La comunicación proactiva ---agenda de
entregas, resúmenes ejecutivos, anticipación de la información--- es aquí más
determinante que en el procedimiento ordinario.

%==============================================================================
\bloqueCronograma
%==============================================================================

\begin{checklist}[Secuencia de trabajo]
  \item \textbf{Semanas $-6$ a $-2$.} Reconstrucción de la contabilidad, verificación
  documental y construcción del mapa de acreedores.
  \item \textbf{Semanas $-2$ a $0$.} Redacción de la propuesta principal \emph{y} de la
  de reintento. Obtención de adhesiones preliminares.
  \item \textbf{Día 0.} Presentación.
  \item \textbf{Resolución de reorganización.} Comienza el período de protección financiera
  simplificado: \destacar{30 días} desde la publicación de la resolución de admisibilidad en el
  \boletin{}, prorrogable \destacar{por una sola vez hasta 15 días} adicionales con el apoyo de al
  menos dos acreedores que representen más del 50\,\% del pasivo.\verificar{Cotejar el plazo base y
  la condición de prórroga contra el articulado vigente de la reorganización simplificada.}
  \item \textbf{Verificación y objeciones.} Control diario del \boletin{}.
  \item \textbf{Junta.} Votación de la propuesta.
  \item \textbf{Rechazo.} Activación inmediata del reintento del artículo 286 S dentro de
  su plazo fatal.
\end{checklist}

%==============================================================================
\bloqueErrores
%==============================================================================

\errorfrecuente{Tratar la declaración jurada como un formalismo}
{Las inexactitudes fundan objeciones e impugnaciones, y en caso de liquidación posterior
alimentan el incidente de mala fe.}
{Construirla con la misma verificación documental que exigiría una auditoría.}

\errorfrecuente{No preparar la propuesta de reintento}
{El plazo del artículo 286 S vence sin poder usarse y el procedimiento deriva en
liquidación simplificada.}
{Redactarla y obtener adhesiones antes de la junta.}

\errorfrecuente{Confundir el patrimonio del socio con el de la empresa}
{Aparecen bienes o deudas cuya titularidad no corresponde al sujeto del concurso.}
{Delimitar ambos patrimonios en el diagnóstico y decidir si se requieren procedimientos
paralelos.}

\errorfrecuente{Presentar cerca del umbral de ventas sin respaldo}
{Un acreedor objeta la procedencia de la vía simplificada y el procedimiento se cae.}
{Acompañar desde el inicio la carpeta tributaria que acredite el umbral con holgura, o
elegir la vía ordinaria.}
